\section{Sources and the literature boundary}\label{sec:sources}
The literature check distinguishes a theorem statement, an exact hypothesis leaf, a source dependent application, and a current reported claim. The accompanying source and leaf manifests identify the versions and inspected locations. Original project reports retain their historical source status, including unsuccessful retrievals and later corrections.

The source basis for the final restricted theorem is Milne's finite set Poitou--Tate duality and local duality \cite{milneADT2006}. Class field conventions, ray reciprocity, restriction with corestriction, and Brauer localization use the authenticated course notes \cite{milneCFT2020}. The principal ideal and unit arguments have been written out in the reader rather than deferred to a computed class number. Marked periods use the stated degeneration and uniformization results \cite{mumford1972,ichikawa1993,gregory-liedtke2019}. Their application to the particular natural family remains an explicit inherited geometry input.

For the wider signature context, Darmon--Granville gives fixed signature finiteness in its declared hyperbolic range \cite{darmon-granville1995}. The Pacetti--Villagra Torcomian preprint supplies the family and specific modular results with hypotheses that must be read individually \cite{pacetti-villagra2025}. In particular, a large prime irreducibility theorem whose cutoff excludes seven is not justified at seven by an unrelated resultant being nonzero modulo seven. The direct project proof has its own arithmetic dependencies.

A current source check also found Chocian's preprint of 22 September 2026, which reports a computer assisted proof for the primitive equation $x^3+y^5=z^7$ \cite{chocian2026}. Its existence and reported scope are recorded here. This release does not independently audit that claimed proof or import it as a theorem closing the recovered project route. Earlier project descriptions of the route as incomplete should therefore be read as descriptions of this project's evidence, not as an exhaustive assertion about all current literature.

The formal source count, exact finite tests, and historical block numbers do not establish novelty. No novelty claim is inferred from a limited source search. The value of this release is a reproducible account of the actual constructions, proved reductions, corrections, and the precisely isolated arithmetic input that the final affine operator still requires.

\clearpage
\begin{thebibliography}{99}\raggedright
\bibitem{mauldin1997} R. D. Mauldin, A generalization of Fermat's Last Theorem: The Beal Conjecture and prize problem, \emph{Notices Amer. Math. Soc.} \textbf{44} (1997), 1436--1437. \url{https://www.ams.org/notices/199711/beal.pdf}.
\bibitem{mauldin1998translation} R. D. Mauldin, \emph{Zobecn\v{e}n\'{i} velk\'{e} Fermatovy v\v{e}ty: Bealova domn\v{e}nka a probl\'{e}m o odm\v{e}n\v{e}}, translated by M. K\v{r}\'{i}\v{z}ek and K. Segeth, \emph{Pokroky matematiky, fyziky a astronomie} \textbf{43} (1998), 104--107. \url{https://dml.cz/handle/10338.dmlcz/137538}.
\bibitem{wiles1995} A. Wiles, Modular elliptic curves and Fermat's Last Theorem, \emph{Ann. of Math.} (2) \textbf{141} (1995), 443--551. DOI: \href{https://doi.org/10.2307/2118559}{10.2307/2118559}.
\bibitem{taylor-wiles1995} R. Taylor and A. Wiles, Ring-theoretic properties of certain Hecke algebras, \emph{Ann. of Math.} (2) \textbf{141} (1995), 553--572. DOI: \href{https://doi.org/10.2307/2118560}{10.2307/2118560}.
\bibitem{mihailescu2004} P. Mih\u{a}ilescu, Primary cyclotomic units and a proof of Catalan's conjecture, \emph{J. reine angew. Math.} \textbf{572} (2004), 167--195. DOI: \href{https://doi.org/10.1515/crll.2004.048}{10.1515/crll.2004.048}.
\bibitem{mihailescu2005} P. Mih\u{a}ilescu, Reflection, Bernoulli numbers and the proof of Catalan's conjecture, in \emph{European Congress of Mathematics}, European Mathematical Society, 2005, 325--340. DOI: \href{https://doi.org/10.4171/009-1/21}{10.4171/009-1/21}.
\bibitem{bilu2005} Y. F. Bilu, Catalan without logarithmic forms (after Bugeaud, Hanrot and Mih\u{a}ilescu), \emph{J. Th\'{e}or. Nombres Bordeaux} \textbf{17} (2005), 69--85. DOI: \href{https://doi.org/10.5802/jtnb.478}{10.5802/jtnb.478}.
\bibitem{darmon-granville1995} H. Darmon and A. Granville, On the equations $z^m=F(x,y)$ and $Ax^p+By^q=Cz^r$, \emph{Bull. London Math. Soc.} \textbf{27} (1995), 513--543. DOI: \href{https://doi.org/10.1112/blms/27.6.513}{10.1112/blms/27.6.513}.
\bibitem{dahmen-siksek2024} S. R. Dahmen and S. Siksek, \emph{On the generalized Fermat equation $x^5+y^3=z^7$}, author working paper, July 2024. \url{https://few.vu.nl/~sdn249/GFE357.pdf}.
\bibitem{pacetti-villagra2025} A. Pacetti and L. Villagra Torcomian, \emph{On the generalized Fermat equation of signature $(5,p,3)$}, arXiv:2512.17845v1, 19 December 2025. \url{https://arxiv.org/abs/2512.17845v1}.
\bibitem{chocian2026} P. Chocian, \emph{The primitive generalized Fermat equation $x^3+y^5=z^7$: A computer-assisted proof}, arXiv:2609.26996v1, 22 September 2026. \url{https://arxiv.org/abs/2609.26996v1}. Recent preprint, not used here as an independently audited theorem.
\bibitem{milneADT2006} J. S. Milne, \emph{Arithmetic Duality Theorems}, second edition, BookSurge, 2006. \url{https://www.jmilne.org/math/Books/ADTnot.pdf}.
\bibitem{milneCFT2020} J. S. Milne, \emph{Class Field Theory}, course notes, version 4.03, 6 August 2020. \url{https://www.jmilne.org/math/CourseNotes/CFT.pdf}.
\bibitem{milneANT2020} J. S. Milne, \emph{Algebraic Number Theory}, course notes, version 3.08, 19 July 2020. \url{https://www.jmilne.org/math/CourseNotes/ANT.pdf}.
\bibitem{milneAV2008} J. S. Milne, \emph{Abelian Varieties}, course notes, version 2.0, March 2008. \url{https://www.jmilne.org/math/CourseNotes/AV.pdf}.
\bibitem{masley1976} J. M. Masley, Solution of small class number problems for cyclotomic fields, \emph{Compositio Math.} \textbf{33} (1976), 179--186. \url{https://www.numdam.org/item/CM_1976__33_2_179_0/}.
\bibitem{mumford1972} D. Mumford, An analytic construction of degenerating curves over complete local rings, \emph{Compositio Math.} \textbf{24} (1972), 129--174. \url{https://www.numdam.org/item/CM_1972__24_2_129_0/}.
\bibitem{ichikawa1993} T. Ichikawa, The universal periods of curves and the Schottky problem, \emph{Compositio Math.} \textbf{85} (1993), 1--8. \url{https://www.numdam.org/item/CM_1993__85_1_1_0/}.
\bibitem{gregory-liedtke2019} O. Gregory and C. Liedtke, $p$-Adic Tate conjectures and abeloid varieties, \emph{Documenta Math.} \textbf{24} (2019), 1879--1934. DOI: \href{https://doi.org/10.4171/DM/718}{10.4171/DM/718}.
\bibitem{baker-rabinoff2013} M. Baker and J. Rabinoff, \emph{The skeleton of the Jacobian, the Jacobian of the skeleton, and lifting meromorphic functions from tropical to algebraic curves}, arXiv:1308.3864v1, 18 August 2013. \url{https://arxiv.org/abs/1308.3864v1}.
\bibitem{dolgachev-lehavi2008} I. Dolgachev and D. Lehavi, \emph{On isogenous principally polarized abelian surfaces}, arXiv:0710.1298v3, 8 March 2008. \url{https://arxiv.org/abs/0710.1298v3}.
\bibitem{ophir-weiss2024} A. Ophir and A. Weiss, \emph{On Ribet's lemma for $\mathrm{GL}_2$ modulo prime powers}, arXiv:2111.01559v3, 1 February 2024. \url{https://arxiv.org/abs/2111.01559v3}.
\bibitem{galbraith2012} S. D. Galbraith, \emph{Mathematics of Public Key Cryptography}, Cambridge University Press, 2012. Chapter 10 checked in the author's extended and corrected version 2.0. \url{https://www.math.auckland.ac.nz/~sgal018/crypto-book/ch10.pdf}.
\bibitem{minac-tan2016} J. Min\'{a}\v{c} and N. D. T\^{a}n, Triple Massey products vanish over all fields, \emph{J. London Math. Soc.} (2) \textbf{94} (2016), 909--932. DOI: \href{https://doi.org/10.1112/jlms/jdw064}{10.1112/jlms/jdw064}. Checked version: arXiv:1412.7611v3, 23 August 2016.
\bibitem{minac-tan2017} J. Min\'{a}\v{c} and N. D. T\^{a}n, Triple Massey products and Galois theory, \emph{J. Eur. Math. Soc.} \textbf{19} (2017), 255--284. DOI: \href{https://doi.org/10.4171/JEMS/665}{10.4171/JEMS/665}.
\bibitem{maire-et-al2026} C. Maire, J. Min\'{a}\v{c}, R. Ramakrishna and N. D. T\^{a}n, On the strong Massey vanishing property for number fields, \emph{Math. Ann.} \textbf{395} (2026), article 102. Published 7 July 2026. DOI: \href{https://doi.org/10.1007/s00208-026-03518-1}{10.1007/s00208-026-03518-1}.
\bibitem{chebolu-minac2009} S. K. Chebolu and J. Min\'{a}\v{c}, Absolute Galois groups viewed from small quotients and the Bloch--Kato conjecture, \emph{Geometry \& Topology Monographs} \textbf{16} (2009), 31--47. DOI: \href{https://doi.org/10.2140/gtm.2009.16.31}{10.2140/gtm.2009.16.31}. The local freeness result is recalled in this survey, not originally proved there.
\bibitem{neukirch1973} J. Neukirch, \"{U}ber das Einbettungsproblem der algebraischen Zahlentheorie, \emph{Invent. Math.} \textbf{21} (1973), 59--116. Korollar 6.4(c), p.90. Digitized published source: \url{https://gdz.sub.uni-goettingen.de/download/pdf/PPN356556735_0021/PPN356556735_0021.pdf}.
\bibitem{neftin2013} D. Neftin, Tamely ramified subfields of division algebras, \emph{J. Algebra} \textbf{378} (2013), 184--195. DOI: \href{https://doi.org/10.1016/j.jalgebra.2012.11.039}{10.1016/j.jalgebra.2012.11.039}. Theorem 2.1 and Corollary 2.3 checked in the author's 16-page manuscript: \url{https://neftin.net.technion.ac.il/files/2015/10/20-6-12.pdf}.
\end{thebibliography}
